\chapter{特殊函数}
\intro{本章介绍贝塞尔函数、勒让德多项式等特殊函数及其在 PDE 中的应用。}

\section{贝塞尔函数}

\begin{mydefinition}{贝塞尔方程}
\[
x^2y''+xy'+(x^2-n^2)y=0
\]
\end{mydefinition}

\begin{theorem}[贝塞尔函数]
第一类贝塞尔函数：
\[
J_n(x)=\sum_{k=0}^{\infty}\frac{(-1)^k}{k!\,\Gamma(k+n+1)}\left(\frac{x}{2}\right)^{2k+n}
\]
\end{theorem}

\begin{figure}[htbp]
\centering
\begin{tikzpicture}
\begin{axis}[
	axis lines=middle,
	xlabel=$x$,
	ylabel={$u(x)$},
	xmin=0,xmax=10,
	ymin=-0.5,ymax=1.2,
	legend pos=north east,
	grid=major,
	width=10cm,height=6cm,
]
\addplot[blue,thick,domain=0:10,samples=200]{cos(deg(x))*exp(-x/8)};
\addplot[red,thick,domain=0:10,samples=200]{sin(deg(x))*exp(-x/8)};
\legend{衰减振荡型解,相移型解}
\end{axis}
\end{tikzpicture}
\caption{贝塞尔函数定性行为}
\end{figure}

\section{勒让德多项式}

\begin{mydefinition}{勒让德方程}
\[
(1-x^2)y''-2xy'+n(n+1)y=0
\]
\end{mydefinition}

\begin{theorem}[勒让德多项式]
\[
P_n(x)=\frac{1}{2^n n!}\frac{d^n}{dx^n}(x^2-1)^n
\]
\end{theorem}

\begin{remark}
贝塞尔函数用于圆柱/圆形区域的分离变量，勒让德多项式用于球对称问题。两者都是 Sturm-Liouville 问题的特征函数。
\end{remark}
